\documentclass[../../main.tex]{subfiles}
\begin{document}

\chapter{AI agents and multi-agent architectures in software engineering}
\label{ch:ai_agents_multiagent}

\section{Chapter overview and learning outcomes}
This chapter moves beyond single-turn prompt-response interactions to study autonomous software engineering agents. We examine how foundation models can interact with software environments through tools, plan complex multi-step workflows, self-reflect on compilation errors, and collaborate in specialized multi-agent teams.

Upon completing this chapter, you will be able to:
\begin{itemize}
    \item Formalize the ReAct (Reasoning + Acting) loop \citep{yao2022react} for software development tasks.
    \item Architect tool-use protocols enabling LLMs to interact with Language Server Protocols (LSPs), test runners, and version control (Git).
    \item Explain the Reflexion feedback loop for autonomous code debugging.
    \item Design role-based multi-agent software engineering teams (Product Manager, Architect, Developer, Reviewer, QA).
    \item Formulate memory hierarchies for software agents: working memory, episodic trace logs, and semantic vector stores.
    \item Understand the structure and difficulty of real-world benchmarks such as SWE-bench \citep{jimenez2024swebench}.
\end{itemize}

\section{The ReAct paradigm: thought, action, and observation}
Autonomous agents transition from pure text generation to an iterative action-perception loop in an environment $\mathcal{E}$.

\begin{definition}{ReAct execution loop}{react_loop}
At step $t$, given task goal $G$ and history trace $H_{t-1}$, the agent executes three steps \citep{yao2022react}:
\begin{enumerate}
    \item \textbf{Thought ($c_t$)}: Internal reasoning step generated by the language model:
    \begin{equation}
    c_t \sim \Prob_{\text{LM}}(\cdot \mid G, H_{t-1}).
    \end{equation}
    \item \textbf{Action ($a_t$)}: Structured command selecting a tool and arguments:
    \begin{equation}
    a_t = \text{tool\_name}(\text{args}) \sim \Prob_{\text{LM}}(\cdot \mid G, H_{t-1}, c_t).
    \end{equation}
    \item \textbf{Observation ($o_t$)}: Deterministic feedback returned by the execution environment:
    \begin{equation}
    o_t = \mathcal{E}(a_t), \quad H_t = H_{t-1} \cup (c_t, a_t, o_t).
    \end{equation}
\end{enumerate}
\end{definition}

\begin{intuition}[ReAct as closed-loop feedback control]
A passive LLM acts as an open-loop controller: it outputs an entire 100-line code patch in one shot without knowing whether the code compiles. A ReAct agent functions as a closed-loop feedback controller: it writes a test, runs the test runner, observes compiler errors, navigates to the definition via LSP, and iteratively refines its solution until all tests pass.
\end{intuition}

\section{Multi-agent software development teams}
Single agents often suffer from context overload and role drift when managing entire projects. Multi-agent systems decompose the software engineering process into specialized personas communicating via message passing:
\begin{itemize}
    \item \textbf{Product Manager}: Clarifies requirements and establishes acceptance criteria.
    \item \textbf{Software Architect}: Designs module boundaries, data structures, and APIs.
    \item \textbf{Programmer}: Writes clean code strictly conforming to architectural interfaces.
    \item \textbf{Code Reviewer}: Audits code against security guidelines and coding standards.
    \item \textbf{QA Engineer}: Generates edge-case unit tests and runs verification suites.
\end{itemize}

\begin{examinsight}[Exam trap: agent loops and halting conditions]
A frequent exam question explores failure modes in autonomous SE agent loops.
Without strict halting conditions (e.g., maximum recursion depth, monotonic error decrease, loop detection in tool calls), agents can easily enter infinite cycles—repeatedly fixing one test while breaking another. Production-grade SE agents must incorporate loop-breaking heuristics and budget constraints.
\end{examinsight}

\begin{deepdive}[SWE-bench: the gold standard for autonomous SE agents]
Jimenez et al. \citep{jimenez2024swebench} introduced \textbf{SWE-bench}, a rigorous benchmark of 2,294 real-world GitHub issues extracted from prominent Python repositories (e.g., Django, SymPy, Scikit-learn). Given an issue description, the agent must inspect the repository, locate the bug across thousands of files, write a patch, and pass hidden unit tests without human intervention.
\end{deepdive}

\end{document}
